% This must be in the first 5 lines to tell arXiv to use pdfLaTeX, which is strongly recommended.
\pdfoutput=1
% In particular, the hyperref package requires pdfLaTeX in order to break URLs across lines.

\documentclass[11pt]{article}

% Remove the "review" option to generate the final version.
\usepackage[]{acl}

% Standard package includes
\usepackage{times}
\usepackage{latexsym}

% For proper rendering and hyphenation of words containing Latin characters (including in bib files)
\usepackage[T1]{fontenc}
% For Vietnamese characters
% \usepackage[T5]{fontenc}
% See https://www.latex-project.org/help/documentation/encguide.pdf for other character sets

% This assumes your files are encoded as UTF8
\usepackage[utf8]{inputenc}

% This is not strictly necessary, and may be commented out,
% but it will improve the layout of the manuscript,
% and will typically save some space.
\usepackage{microtype}

% If the title and author information does not fit in the area allocated, uncomment the following
%
%\setlength\titlebox{<dim>}
%
% and set <dim> to something 5cm or larger.

\title{Chart-Aware Transformers: Learning Structured Embeddings for BI Report Understanding}

% Author information can be set in various styles:
% For several authors from the same institution:
% \author{Author 1 \and ... \and Author n \\
%         Address line \\ ... \\ Address line}
% if the names do not fit well on one line use
%         Author 1 \\ {\bf Author 2} \\ ... \\ {\bf Author n} \\
% For authors from different institutions:
% \author{Author 1 \\ Address line \\  ... \\ Address line
%         \And  ... \And
%         Author n \\ Address line \\ ... \\ Address line}
% To start a seperate ``row'' of authors use \AND, as in
% \author{Author 1 \\ Address line \\  ... \\ Address line
%         \AND
%         Author 2 \\ Address line \\ ... \\ Address line \And
%         Author 3 \\ Address line \\ ... \\ Address line}

\author{%
\makebox[\textwidth][c]{%
\begin{tabular}{ccc}
\large Akhil Srivatsava & \large Bimal Dev & \large Prashant Agarwal \\
\normalsize \texttt{sriakhil25@gmail.com} & \normalsize \texttt{bimal.d3v@gmail.com} & \normalsize \texttt{prashant.agarwal92@gmail.com} \\
[0.7em]
\multicolumn{3}{c}{%
  \large Pavan Karumuru\hspace{2em}%
  \large Krishna Rao Kakkirala}\\
\multicolumn{3}{c}{%
  \normalsize \texttt{pavankumar.ds20@gmail.com}\hspace{2em}%
  \normalsize \texttt{krishna.kakkirala@gmail.com}}%
\end{tabular}%
}%
}

\begin{document}
\maketitle
\begin{abstract}
Traditional vision models that tokenize images into uniform patches often fail to capture the structured, semantically rich nature of charts and BI (Business Intelligence) reports. In this work, we introduce Chart-Aware Transformers, a novel architecture that models charts as sequences of structured elements—such as titles, axes, legends, bars, and annotations—rather than pixel-based patches. Each element is embedded with a combination of class, textual, and layout features, enabling the model to encode semantic and spatial relationships critical for downstream tasks like chart retrieval, question answering, and document understanding. By treating charts as compositional units with explicit structure, our approach yields semantically faithful and information-rich representations, outperforming patch-based and purely textual baselines. We demonstrate that this element-centric formulation is better aligned with the nature of chart data, unlocking new capabilities for intelligent document understanding in enterprise and analytics workflows.
\end{abstract}

\section{Introduction}

Charts and visualizations serve as fundamental communication tools across diverse domains including business intelligence, financial reporting, scientific research, and data journalism. In enterprise settings, analysts spend significant time manually interpreting dashboards and reports to extract insights, while accessibility remains a critical challenge for visually impaired users who rely on alternative representations of visual data. Recent advances in chart question answering~\cite{masry2022chartqa} and visual document understanding have demonstrated the potential for automated systems to assist with chart interpretation tasks. However, despite progress in general vision-language models, the structured and compositional nature of charts—with their explicit semantic elements like axes, legends, and data marks—remains underexploited in current approaches, motivating the need for specialized architectures that can reason over chart structure.

Recent work has made significant progress in chart understanding through specialized architectures. UniChart~\cite{masry2023unichart} introduced a chart-focused encoder-decoder framework trained on 611K charts, supporting summarization and QA tasks, though it struggles with compositional reasoning over complex chart structures. ChartAssistant~\cite{zhang2024chartassistant} improved performance through chart-to-table pretraining and multitask instruction tuning, demonstrating strong generalization across chart types. LayoutLMv3~\cite{huang2023layoutlmv3} showed the value of layout-aware modeling for structured documents by integrating text, vision, and layout features. While these approaches have advanced the field, they primarily operate on patch-based image representations or rely on intermediate table conversions, rather than directly modeling the compositional structure of chart elements. This motivates our approach of treating chart components—such as bars, axes, and legends—as explicit structured tokens.

Understanding charts and Business Intelligence (BI) reports is a critical challenge in visual document analysis. Traditional Vision Transformer (ViT)-based models rely on patch-level image tokenization, which often fails to capture the structured, semantic, and spatial relationships between elements in a chart. These models lack an explicit representation of chart components like axes, legends, bars, titles, and annotations — elements essential for interpreting the information encoded in structured visualizations.

In this paper, we introduce Chart-Aware Transformers, a novel architecture that explicitly models chart elements as discrete tokens. Each token is enriched with textual, semantic, and layout-aware embeddings to reflect its class (e.g., “bar”, “x-axis”, “title”), content (e.g., OCR-extracted text), and spatial location (e.g., bounding box coordinates and size). By treating chart elements as the core units of representation — rather than raw pixels — we enable Transformer-based reasoning that aligns more naturally with the structure and semantics of BI charts.

We train and evaluate our system on a large-scale synthetic dataset covering 43 distinct chart element classes with COCO-style annotations. Our element-centric representation preserves semantic and spatial relationships between chart components, enabling more effective reasoning for downstream tasks such as element counting, chart retrieval, and question answering compared to patch-based approaches.

\section{Prior Literature}

Chart understanding tasks require models to reason over structured visual elements, including axes, legends, and data marks, which general-purpose MLLMs like GPT-4V often struggle with due to weak visual-semantic grounding and poor layout sensitivity \cite{li2024visualizationliteracy, huang2024survey}.

Recent specialized models address these limitations. \textbf{UniChart} \cite{masry2023unichart} introduces a chart-focused encoder-decoder framework trained on 611K charts, supporting tasks like summarization and chart QA. While efficient, it struggles with compositional reasoning. \textbf{ChartAssistant} \cite{zhang2024chartassistant} improves performance via a two-stage pipeline (chart-to-table pretraining + multitask instruction tuning) and shows strong generalization across chart types, especially in numerical and referring QA. \textbf{ChartInstruct} \cite{lin2024chartinstruct} complements these with a broad instruction-tuned dataset, enhancing chart reasoning without architectural changes.

To better encode layout-sensitive input, \textbf{LayoutLMv3} \cite{huang2023layoutlmv3} integrates text, vision, and layout using a unified ViT backbone. This layout-aware modeling is better suited to structured documents like charts, especially compared to models like \textbf{MM-GEM} \cite{zhang2024mmgem}, which lack fine-grained spatial grounding.

Our approach builds on these insights by modeling chart elements as structured tokens, enriched with class, text, and layout embeddings. We evaluate on the \textbf{ChartAssistant dataset} \cite{chartassistant_dataset}, which offers realistic chart-QA pairs, supporting fine-grained benchmarking of reasoning over structured visual input.



\section{Data}
To support the training of our chart element detection model, we developed a Python-based chart generation pipeline capable of producing synthetic chart images (including bar, line, pie, and stacked bar charts) with detailed annotations in the COCO format. The pipeline uses libraries such as Matplotlib, NumPy, and Plotly to programmatically generate diverse chart types along with structured annotations for each visual element.

The system defines a comprehensive taxonomy of chart components (e.g., title, axes, axis titles, tick labels, gridlines, legends, bars, lines, and pie slices) and extracts precise bounding boxes for each element. These annotations follow the COCO dataset structure, capturing metadata such as bounding box coordinates, area, category IDs, and image associations.

For each chart type, the pipeline computes spatial annotations using geometric primitives or text layout analysis, ensuring pixel-accurate labels. It supports various configurations, such as line style (e.g., linear, sine, noisy), bar orientation (horizontal or vertical), and pie slice randomization. The resulting dataset includes both the rendered chart images and a corresponding JSON file of annotations, enabling downstream tasks such as chart element detection and visual data extraction.

This synthetic dataset complements our manually annotated and publicly sourced data (e.g., ChartAssistant), providing a rich and customizable training resource for layout-aware chart understanding.

\section{Model}

We design a Transformer-based encoder model to learn structured visual-textual representations from chart annotations. Each chart is represented as a variable-length sequence of elements (e.g., bars, titles, axes), where each element is encoded into a fixed-size embedding using a combination of class, bounding box, and area features.

\subsection*{Embedding Layer}

Each chart element is represented by:
\begin{itemize}
    \item A one-hot encoded \textbf{class vector} (e.g., bar, title, x-axis, y-axis)
    \item A normalized \textbf{bounding box}: $(x, y, w, h)$
    \item A normalized \textbf{area scalar}
\end{itemize}

We use separate Multi-Layer Perceptrons (MLPs) to project each of these features into a shared embedding space of dimension $D=256$, and concatenate the resulting vectors to form the final input embedding of dimension $3D=768$ per token.

\subsection*{Transformer Encoder}

The chart-level sequence of embedded elements is fed into a Transformer encoder consisting of:
\begin{itemize}
    \item Multi-Head Self Attention (MHSA)
    \item Feedforward layers with ReLU activation
    \item LayerNorm and residual connections
    \item Positional encoding via order-preserving padding
\end{itemize}

Unlike traditional autoregressive settings, our self-attention blocks do not use causal masks, as the model observes all elements simultaneously.

\subsection*{Classification Head}

We cast the chart understanding task as a \textbf{counting problem}, specifically predicting the number of \texttt{bar} elements in a chart. This is framed as a multi-class classification task over the range $[0, 100]$. After encoding the sequence, we apply mean pooling over token embeddings and pass the result through a classification head:

\[
\hat{y} = \text{Softmax}(W \cdot \text{MeanPool}(Z) + b)
\]

where $Z$ is the output of the final transformer layer.

\subsection*{Loss Function}

We use cross-entropy loss to train the model to predict the number of bars in each chart. This formulation encourages the model to learn rich token-level representations aligned with global structural patterns.


\section{Methods}

We propose a \textbf{Chart Transformer}, a layout-aware, Transformer-based architecture designed for structured understanding of chart images. The model operates over pre-extracted chart elements—such as bars, titles, axes, and legends—represented using COCO-style annotations. Each element is embedded using its class label, normalized bounding box coordinates, and area. These features are projected into a unified embedding space, enriched with positional encodings, and passed through a Transformer encoder. The self-attention mechanism enables the model to reason over spatial and semantic relationships between chart elements. Chart Transformer supports multiple downstream tasks, including element counting, chart-to-JSON conversion, and semantic type classification.

\subsection*{Dataset Generation}

To train and evaluate our system, we generate a large-scale synthetic dataset using Matplotlib, covering 43 distinct chart element classes including components from bar, line, and area charts. For each chart, we programmatically extract bounding boxes using the Matplotlib rendering backend, capturing pixel-accurate positions of visual elements such as bars, titles, axes, and tick marks. Each chart is annotated in COCO-style JSON format, ensuring compatibility with common vision datasets and workflows. This synthetic dataset complements publicly sourced data and provides a rich, customizable training resource for layout-aware chart understanding.

\subsection*{Feature Representation}

The extracted chart elements are converted into fixed-length feature vectors combining three key components: (1) class embeddings that encode the semantic type of each element (e.g., bar, axis, legend), (2) OCR/text embeddings that capture textual content when present, and (3) normalized layout information including bounding box coordinates and area. These feature vectors are projected into a unified embedding space and fed into the Transformer encoder, enabling the model to reason over both semantic and spatial relationships between chart elements.

\subsection*{Advantages of Element-Centric Representation}

This element-centric formulation offers several key advantages over patch-based approaches:
\begin{itemize}
    \item It preserves semantic and spatial relationships that are lost when charts are tokenized into uniform patches
    \item It enables explicit reasoning over individual elements and their structural roles within the chart
    \item It supports generalization to diverse downstream tasks such as chart-based question answering, retrieval, and summarization
    \item It aligns naturally with the compositional structure of charts, where meaning emerges from relationships between discrete visual components
\end{itemize}

\subsection*{Baseline}

To evaluate the effectiveness of Chart Transformer, we compare it against a \textbf{Vision Transformer (ViT)} baseline. In this setup, the raw chart image is passed through a pre-trained ViT model. The output embedding—typically the \texttt{[CLS]} token—is used by a task-specific prediction head. This baseline lacks explicit structural input and serves as a contrast to our layout-aware Chart Transformer, highlighting the benefit of reasoning over element-level representations.

\subsection*{Evaluation}

We benchmark both models on tasks such as predicting the number of chart elements. Since the task is inherently classification-based, \textbf{accuracy} is used as the primary evaluation metric. This comparison demonstrates the advantage of layout-aware modeling for chart understanding, particularly in classification tasks involving structured visual elements.


\section{Results}

In our experiments, the \textbf{Chart Transformer} was compared against the \textbf{Vision Transformer (ViT)} for the task of predicting the number of chart elements from images. Both models were evaluated using \textbf{accuracy} as the primary metric, since the task is framed as a classification problem.

\subsection*{Performance Comparison}

\begin{itemize}
    \item \textbf{Chart Transformer}: Achieved an accuracy of \textbf{92\%} in predicting the number of chart elements across the evaluation dataset.
    \item \textbf{Vision Transformer (ViT)}: Was consistently outperformed by the Chart Transformer, indicating the benefits of incorporating explicit layout-aware element representations.
\end{itemize}


\section{Analysis} 

The \textbf{Chart Transformer} outperformed the \textbf{Vision Transformer (ViT)}, highlighting the benefits of incorporating layout-aware reasoning through explicit element-level embeddings. While the ViT model is effective for generic vision tasks, it struggled to capture the spatial and semantic relationships between chart components. This limitation led to lower accuracy on tasks that required structural understanding of charts.

In contrast, the Chart Transformer—by leveraging structured element representations and positional encodings—was better suited to model inter-element dependencies. This architectural advantage contributed to improved performance in the classification task, demonstrating the importance of layout sensitivity in chart understanding models.


\section{Conclusion} 
The \textbf{Chart Transformer} demonstrated superior performance compared to the \textbf{Vision Transformer (ViT)} in the task of predicting the number of chart elements, achieving an accuracy of \textbf{92\%}. By utilizing layout-aware embeddings and a custom architecture designed specifically for structured chart data, the Chart Transformer was able to more effectively capture and classify chart elements than the ViT, which lacks such specialized handling.

These results highlight the importance of domain-specific architectures for tasks involving complex visual structures such as charts. The performance of the Chart Transformer reinforces the idea that incorporating knowledge of layout and inter-element relationships can significantly enhance model accuracy in chart understanding tasks.

\subsection*{Future Scope}

Building on the success of the Chart Transformer, future work will explore integrating it with existing \textbf{text encoder models} and the \textbf{Vision Transformer (ViT)} to build a unified system. The aim is to enable retrieval-augmented generation (RAG) pipelines that can jointly process chart visuals and textual descriptions. Such a system would support natural language queries for retrieving and interpreting chart data, advancing the applicability of chart understanding in real-world information retrieval and analytics settings.


\section{Known Project Limitations}

Our work is subject to several limitations that may impact its generalization and real-world applicability.

First, our training data is largely synthetic and may not reflect the full visual diversity of real-world charts, leading to possible performance drops in more complex or noisy settings. Second, the annotation schema prioritizes structural elements but may miss finer visual or stylistic details.

Additionally, our task framing simplifies chart understanding to classification and counting, which does not capture the full spectrum of reasoning often required in practical BI applications. Evaluation is limited to English-language datasets and primarily Western-style chart formats, which may introduce cultural or linguistic biases.

Users applying this system in diverse domains or languages should account for these constraints, and future work will aim to improve generalization, reasoning depth, and multilingual support.

% Entries for the entire Anthology, followed by custom entries
\bibliography{anthology,custom}

\appendix

\end{document}
